\chapter*{Acknowledgments／致谢}
\addcontentsline{toc}{chapter}{Acknowledgments／致谢}
\EN{This work was supported by the National Natural Science Foundation of China (Grant Nos. 12471234, 12171439) and Science Foundation of Henan Academy of Sciences (Grant No. 20252319002).}
\CN{本工作得到国家自然科学基金（项目号 12471234、12171439）及河南省科学院科学基金（项目号 20252319002）支持。}
\section*{Data Availability Statement／数据可用性声明}
\EN{Data sharing not applicable to this article as no data sets were generated or analyzed during the current study.}
\CN{本研究未生成或分析数据集，因此本文不适用数据共享声明。}
\chapter*{References／参考文献}
\addcontentsline{toc}{chapter}{References／参考文献}
\begin{note}
以下保留原文 33 个参考文献序号及著录，题名附中文译名。外部文献的历史结论按原作者引用理解，未核读的文献不作为本笔记独立验证的来源。
原文 [1]、[2] 的作者名印为 Hirita，[3] 印为 Green，照录保留；常见规范拼写分别为 Hirota、Greene。
本文期刊原件：H. Zhang, X. Geng, H. Liu, \emph{Studies in Applied Mathematics} 155 (2025), e70113，\href{https://doi.org/10.1111/sapm.70113}{doi:10.1111/sapm.70113}。
\end{note}
\begin{enumerate}[label={[\arabic*]},leftmargin=2.8em]
\item R. Hirita and J. Satsuma, ``Soliton Solutions of a Coupled Korteweg-de Vries Equation,'' \emph{Physics Letters A} 85 (1981): 407--408.
\CN{耦合 Korteweg--de Vries 方程的孤子解。}
\item R. Hirita and J. Satsuma, ``A Coupled KdV Equation is One Case of the Four-Reduction of the KP Hierarchy,'' \emph{Journal of the Physical Society of Japan} 51 (1982): 3390--3397.
\CN{一种耦合 KdV 方程是 KP 层级四约化的一个特例。}
\item C. S. Gardner, J. M. Green, M. D. Kruskal, and R. M. Miura, ``Method for Solving the Korteweg--de Vries Equation,'' \emph{Physical Review Letters} 19 (1967): 1095--1097.
\CN{求解 Korteweg--de Vries 方程的方法。}
\item K. Grunert and G. Teschl, ``Long-Time Asymptotics for the Korteweg--de Vries Equation via Nonlinear Steepest Descent,'' \emph{Mathematical Physics, Analysis and Geometry} 12 (2009): 287--324.
\CN{通过非线性最速下降法研究 KdV 方程的长时间渐近行为。}
\item R. Dodd and A. P. Fordy, ``On the Integrability of a System of a Coupled KdV Equation,'' \emph{Physics Letters A} 89 (1982): 168--170.
\CN{论一个耦合 KdV 方程组的可积性。}
\item W. H. Tam, W. X. Ma, X. B. Hu, and D. L. Wang, ``The Hirota--Satsuma Coupled KdV Equation and a Coupled Ito System Revisited,'' \emph{Journal of the Physical Society of Japan} 69 (2000): 45--52.
\CN{再论 Hirota--Satsuma 耦合 KdV 方程与耦合 Ito 系统。}
\item D. Levi, ``A Hierarchy of Coupled Korteweg--de Vries Equation,'' \emph{Physics Letters A} 95 (1983): 7--10.
\CN{一个耦合 KdV 方程层级。}
\item H. C. Hu and Q. P. Liu, ``New Darboux Transformation for Hirota--Satsuma Coupled KdV System,'' \emph{Chaos, Solitons \& Fractals} 17 (2003): 921--928.
\CN{Hirota--Satsuma 耦合 KdV 系统的新 Darboux 变换。}
\item\phantomsection\label{ref:original9} Y. T. Wu, X. G. Geng, X. B. Hu, and S. M. Zhu, ``A Generalized Hirota--Satsuma Coupled Korteweg--de Vries Equation and Miura Transformations,'' \emph{Physics Letters A} 255 (1999): 259--264.
\CN{广义 Hirota--Satsuma 耦合 KdV 方程及 Miura 变换。}
\item X. Zhou, ``The Riemann-Hilbert Problem and Inverse Scattering,'' \emph{SIAM Journal on Mathematical Analysis} 20 (1989): 966--986.
\CN{Riemann--Hilbert 问题与逆散射。}
\item P. Deift and X. Zhou, ``A Steepest Descent Method for Oscillatory Riemann--Hilbert Problems. Asymptotics for the MKdV Equation,'' \emph{Annals of Mathematics} 137 (1993): 295--368.
\CN{振荡 RH 问题的最速下降法及 mKdV 方程的渐近行为。}
\item X. G. Geng and H. Liu, ``The Nonlinear Steepest Descent Method to Long-Time Asymptotics of the Coupled Nonlinear Schrödinger Equation,'' \emph{Journal of Nonlinear Science} 28 (2018): 739--763.
\CN{耦合非线性薛定谔方程长时间渐近的非线性最速下降法。}
\item X. G. Geng, K. D. Wang, and M. M. Chen, ``Long-Time Asymptotics for the Spin-1 Gross--Pitaevskii Equation,'' \emph{Communications in Mathematical Physics} 382 (2021): 585--611.
\CN{自旋为一的 Gross--Pitaevskii 方程的长时间渐近。}
\item X. G. Geng, J. Wang, K. D. Wang, and R. M. Li, ``Soliton Resolution for the Complex Short-Pulse Positive Flow With Weighted Sobolev Initial Data in the Space-Time Soliton Regions,'' \emph{Journal of Differential Equations} 386 (2024): 214--268.
\CN{加权 Sobolev 初值的复短脉冲正流在时空孤子区中的孤子分解。}
\item R. Beals and R. R. Coifman, ``Scattering and Inverse Scattering for First Order Systems,'' \emph{Communications on Pure and Applied Mathematics} 37 (1984): 39--90.
\CN{一阶系统的散射与逆散射。}
\item R. Beals, P. Deift, and C. Tomei, \emph{Direct and Inverse Scattering on the Line} (American Mathematical Society, 1988).
\CN{直线上的正散射与逆散射。}
\item B. Prinari, G. Biondini, and A. D. Trubatch, ``Inverse Scattering Transform for the Multi-Component Nonlinear Schrödinger Equation With Nonzero Boundary Conditions,'' \emph{Studies in Applied Mathematics} 126 (2011): 245--302.
\CN{非零边界条件下多分量非线性薛定谔方程的逆散射变换。}
\item G. Biondini, D. Kraus, and B. Prinari, ``The Three-Component Defocusing Nonlinear Schrödinger Equation With Nonzero Boundary Conditions,'' \emph{Communications in Mathematical Physics} 348 (2016): 475--533.
\CN{非零边界条件下三分量散焦非线性薛定谔方程。}
\item A. B. de Monvel and D. Shepelsky, ``A Riemann-Hilbert Approach for the Degasperis-Procesi Equation,'' \emph{Nonlinearity} 26 (2013): 2081--2107.
\CN{Degasperis--Procesi 方程的 RH 方法。}
\item A. B. de Monvel, D. Shepelsky, and L. Zielinski, ``A Riemann--Hilbert Approach for the Novikov Equation,'' \emph{Symmetry, Integrability and Geometry: Methods and Applications} 12 (2016): 1--22.
\CN{Novikov 方程的 RH 方法。}
\item A. B. de Monvel and D. Shepelsky, ``The Ostrovsky-Vakhnenko Equation by a Riemann-Hilbert Approach,'' \emph{Journal of Physics A} 48 (2015): 035204.
\CN{用 RH 方法研究 Ostrovsky--Vakhnenko 方程。}
\item C. Charlier and J. Lenells, ``The `Good' Boussinesq Equation: A Riemann--Hilbert Approach,'' \emph{Indiana University Mathematics Journal} 71 (2022): 1505--1562.
\CN{良性 Boussinesq 方程的 RH 方法。}
\item C. Charlier, J. Lenells, and D. S. Wang, ``The `Good' Boussinesq Equation: Long-Time Asymptotics,'' \emph{Analysis \& PDE} 6 (2023): 1351--1388.
\CN{良性 Boussinesq 方程的长时间渐近。卷号按原文照录。}
\item C. Charlier and J. Lenells, ``On Boussinesq's Equation for Water Waves,'' preprint, 2023, doi:10.48550/arXiv.2204.02365 [nlin.SI].
\CN{论水波的 Boussinesq 方程。}
\item D. S. Wang and X. D. Zhu, ``Long-Time Asymptotics of the Sawada-Kotera Equation and Kaup-Kupershmidt Equation on the Line,'' preprint, 2023, doi:10.48550/arXiv.2307.08196 [nlin.SI].
\CN{直线上 Sawada--Kotera 与 Kaup--Kupershmidt 方程的长时间渐近。}
\item D. S. Wang, C. Zhu, and X. D. Zhu, ``Miura Transformations and Large-Time Behaviors of the Hirota-Satsuma Equation,'' preprint, 2024, doi:10.48550/arXiv.2404.01215 [nlin.SI].
\CN{Hirota--Satsuma 方程的 Miura 变换与大时间行为。原文 DOI 域名中的排印点号在此改正。}
\item L. Huang, D. S. Wang, and X. D. Zhu, ``Long-Time Asymptotics of the Tzitzéica Equation on the Line,'' preprint, 2024, doi:10.48550/arXiv.2404.04999 [nlin.SI].
\CN{直线上 Tzitzéica 方程的长时间渐近。}
\item P. J. Caudrey, ``The Inverse Problem for the Third Order Equation $u_{xxx}+q(x)u_x+r(x)u=-\ii\zeta^3u$,'' \emph{Physics Letters A} 79 (1980): 1819--1826.
\CN{三阶谱方程的逆问题。}
\item P. Deift and C. Tomei, ``Inverse Scattering and the Boussinesq Equation,'' \emph{Communications on Pure and Applied Mathematics} 35 (1982): 567--628.
\CN{逆散射与 Boussinesq 方程。}
\item L. Huang and J. Lenells, ``Nonlinear Fourier Transforms for the Sine-Gordon Equation in the Quarter Plane,'' \emph{Journal of Differential Equations} 264 (2018): 3445--3499.
\CN{四分之一平面上 sine--Gordon 方程的非线性 Fourier 变换。}
\item H. Liu, J. Shen, and X. G. Geng, ``Inverse Scattering Transformation for the N-Component Focusing Nonlinear Schrödinger Equation With Nonzero Boundary Conditions,'' \emph{Letters in Mathematical Physics} 113 (2023): 1--47.
\CN{非零边界条件下 N 分量聚焦非线性薛定谔方程的逆散射变换。}
\item P. J. Caudrey, ``The Inverse Problem for a General $N\times N$ Spectral Equation,'' \emph{Physica D} 6 (1982): 51--66.
\CN{一般 N 阶矩阵谱方程的逆问题。}
\item V. E. Zakharov and A. B. Shabat, ``A Plan for Integrating the Nonlinear Equations of Mathematical Physics by the Method of the Inverse Scattering Problem. I,'' \emph{Functional Analysis and Its Applications} 8 (1974): 226--235.
\CN{用逆散射问题方法积分数学物理非线性方程的方案，第一部分。}
\end{enumerate}
\section*{Additional background／补充先修资料}
范恩贵：《可积系统、正交多项式和随机矩阵——Riemann--Hilbert 方法》，仓库中的二校版（文件日期标识 20220303）。以其前六章作为阅读背景；新用到的外幂与 Fredholm 构造在正文实际使用处解释。
\section*{Publication notice／原件出版说明}
\EN{\copyright\ 2025 Wiley Periodicals LLC.}
\CN{原件出版版权：2025 年 Wiley Periodicals LLC。}
\EN{Copyright of Studies in Applied Mathematics is the property of Wiley-Blackwell and its content may not be copied or emailed to multiple sites or posted to a listserv without the copyright holder's express written permission. However, users may print, download, or email articles for individual use.}
\CN{原件附页说明：《Studies in Applied Mathematics》的版权属于 Wiley--Blackwell；未经版权人明确书面许可，其内容不得复制或群发至多个站点或邮件列表。用户可以为个人使用打印、下载或发送文章。}
